\section[$q$-difference operators]{$\boldsymbol{q}$-difference operators} \label{q-dif}

In this section we introduce notation and recall basic facts about algebra $\qD$, see \cite{FFZ, GL,KR93}.
\begin{Definition}
The associative algebra of $q$-difference operators $\qaD$ is an associative algebra generated by $D^{\pm 1}$ and $x^{\pm 1}$ with the relation $Dx=qxD$.
\end{Definition}

\begin{Definition}
The algebra of $q$-difference operators $\qD$ is a Lie algebra with a basis $E_{k,l}$ {\upshape(}where $(k, l) \in \mathbb{Z}^2 \backslash \{(0,0) \}${\upshape)}, $c$ and $c'$. The elements $c$ and $c'$ are central. All other commutators are given by
\begin{gather}
[ E_{k,l} , E_{r,s} ] = \big( q^{(sk-lr)/2}-q^{(lr-sk)/2} \big) E_{k+r, l+s} + \delta_{k, -r} \, \delta_{l, -s} ( c' k + c l ). \label{qDc}
\end{gather}
\end{Definition}

\begin{Remark}Note that the vector subspace of $\qaD$ spanned by $x^l D^k$ (for $(l,k) \neq (0,0)$) is closed under commutation, i.e., has a natural structure of Lie algebra (denote this Lie algebra by $\mathrm{Diff}_q^{\mathrm{L}}$). Consider a basis of this Lie algebra $E_{k,l} := q^{kl/2} x^l D^k$. Finally, $\qD$ is a central extension of $\mathrm{Diff}_q^{\mathrm{L}}$ by two-dimensional abelian Lie algebra spanned by $c$ and $c'$.
\end{Remark}

\subsection[${\rm SL}_2(\mathbb{Z} )$ action]{$\boldsymbol{{\rm SL}_2(\mathbb{Z} )}$ action}

In this section we will define action ${\rm SL}_2 ( \mathbb{Z} )$ on $\qD$. Let $\sigma$ be an element of ${\rm SL}_2 ( \mathbb{Z} )$ corresponding to a matrix
\begin{gather*}
\sigma = \begin{pmatrix}
m' & m \\
n' & n
\end{pmatrix}. %\label{notation:sigma}
\end{gather*}
Then $\sigma$ acts as follows
\begin{gather} \label{action1}
\sigma( E_{k,l} ) = E_{ m' k + m l, n' k + n l }, \qquad
\sigma (c') = m' c' + n' c, \qquad \sigma (c) = m c' + n c.
\end{gather}

\begin{Proposition} Formula \eqref{action1} defines ${\rm SL}_2 ( \mathbb{Z} )$ action on $\qD$ by Lie algebra automorphisms.
\end{Proposition}
\begin{proof} Note that \eqref{qDc} is ${\rm SL}_2 ( \mathbb{Z} )$ covariant. \end{proof}

For any $\qD$-module $M$ denote by $\rho_{M} \colon \qD \rightarrow \mathfrak{gl} (M)$ the corresponding homomorphism.

\begin{Definition} \label{twisted} For any $\qD$-module $M$ and $\sigma \in {\rm SL}(2,\mathbb{Z})$ let us define the representation $M^\sigma$ as follows. $M$ and $M^{\sigma}$ are the same vector space with different actions, namely $\rho_{M^{\sigma}} = \rho_M \circ \sigma$.
\end{Definition}

We will refer to $M^{\sigma}$ as a \emph{twisted representation}. More precisely, $M^{\sigma}$ is the representation $M$, twisted by~$\sigma$.

\subsection{Chevalley generators and relations}

The Lie algebra $\qD$ is generated by $E_k := E_{1, k}$, $F_k:=E_{-1, k}$ and $H_k := E_{0, k}$. We will call them the \emph{Chevalley generators} of $\qD$.
Define the following \emph{currents} (i.e., formal power series with coefficients in $\qD$)
\begin{gather*}
E(z) = \sum_{k \in \mathbb{Z}} E_{1,k} z^{-k} = \sum_{k \in \mathbb{Z}} E_k z^{-k},\\
F(z) = \sum_{k \in \mathbb{Z}} E_{-1, k} z^{-k} = \sum_{n \in \mathbb{Z}} F_k z^{-k},\\
H(z) = \sum_{k \neq 0} E_{0, k} z^{-k} = \sum_{k \neq 0} H_k z^{-k}.
\end{gather*}

Let us also define the formal delta function
\begin{gather*}
\delta(x) = \sum_{k \in \mathbb{Z}} x^k.
\end{gather*}
\begin{Proposition}\label{relation} Lie algebra $\qD$ is presented by the generators $E_k$, $F_k$ $($for all $k \in \mathbb{Z})$, $H_l$ {\upshape(}for $l \in \mathbb{Z} \backslash \{0 \}${\upshape)}, $c$, $c'$ and the following relations
\begin{gather}
[H_k, H_l] = k c \delta_{k+l, 0}, \label{RqDHH} \\
[H_k , E(z)] = \big(q^{-k/2}-q^{k/2}\big)z^k E(z), \qquad
[H_k, F(z)]= \big(q^{k/2}-q^{-k/2}\big) z^k F(z), \label{RqDHE-F} \\
(z-qw)\big(z-q^{-1} w\big) [E(z), E(w)]= 0, \qquad
(z-qw)\big(z-q^{-1} w\big) [F(z), F(w)]= 0 , \label{RqDE2} \\
[ E (z) , F (w)] = \big( H\big( q^{-1/2} w \big) - H\big(q^{1/2} w\big) + c'\big) \delta( w / z ) + c \frac{w}{z} \delta'(w/ z),\label{RqDEF} \\
z_2 z_3^{-1} [E(z_1), [E(z_2), E(z_3)]] + \mathrm{cyclic} =0, \label{RqDEEE} \\
z_2 z_3^{-1} [F(z_1), [F(z_2), F(z_3)]] + \mathrm{cyclic} =0. \label{RqDFFF}
\end{gather}
\end{Proposition}
One can find a proof of Proposition \ref{relation} in \cite[Theorem~2.1]{M07} or \cite[Theorem~5.5]{Ts}.

\section{Fock module} \label{level1}
In this section we review basic constructions of representations of $\qD$ with $c=1$ and $c'=0$. These construction were studied in~\cite{GL}.

\subsection{Free boson realization} \label{boson}

Introduce the Heisenberg algebra generated by $a_k$ (for $k \in \mathbb{Z}$) with relation $[a_k, a_l] = k \delta_{k+l, 0}$. Consider the Fock module $F_{\alpha}^a$ generated by $| \alpha \rangle$ such that $a_k | \alpha \rangle = 0$ for $k > 0$, $a_0 | \alpha \rangle = \alpha | \alpha \rangle$.
\begin{Proposition} \label{prop:sec3 Boson Fock}
The following formulas
determine an action of $\qD$ on $F_{\alpha}^a$:
\begin{gather}
c \mapsto 1, \qquad c' \mapsto 0, \qquad H_k \mapsto a_k, \label{boson1}\\
E(z) \mapsto \frac{u}{1-q} \exp \left( \sum_{k>0} \frac{q^{-k/2}-q^{k/2}}{k} a_{-k} z^k \right) \exp \left( \sum_{k<0} \frac{q^{-k/2}-q^{k/2}}{k} a_{-k} z^k \right) \label{eq:E},\\
F(z) \mapsto \frac{u^{-1}}{1-q^{-1}} \exp \left( \sum_{k>0} \frac{q^{k/2}-q^{-k/2}}{k} a_{-k} z^k \right) \exp \left( \sum_{k<0}\frac{q^{k/2}-q^{-k/2}}{k} a_{-k} z^k \right). \label{eq:F}
\end{gather}
\end{Proposition}

We will denote this representation by $\mathcal{F}_u$.

\begin{Remark} \label{remark:NOalpha}Note that $\alpha$ does not appear in formulas \eqref{boson1}--\eqref{eq:F}. But we will need operator $a_0$ later (see the proof of Proposition \ref{Boson-Fermion}) for the boson-fermion correspondence. Heuristically, one can think that $u = q^{-\alpha}$.
\end{Remark}

\begin{Remark}[on our notation]In this paper, we consider several algebras and their action on the corresponding Fock modules. We choose the following notation. All these representations are denoted by the letter F (for Fock) with some superscript to mention an algebra. Since $\qD$ is the most important algebra in our paper, we use no superscript for its representation.
Also, let us remark that we consider several copies of the Heisenberg algebra. To distinguish their Fock modules, we write a letter for generators as a superscript.
\end{Remark}

The standard bilinear form on $F^a_{\alpha}$ is defined by the following conditions: operator $a_{-k}$ is dual of $a_k$, the pairing of $\valpha$ with itself equals 1. We will use the bra-ket notation for this scalar product. For an operator $A$ we denote by $\langle \alpha | A | \alpha \rangle$ the scalar product of $A | \alpha \rangle$ with $\valpha$.

\begin{Proposition}\label{Fock} Suppose the algebra $\qD$ acts on $F^a_{\alpha}$ so that $H_k \mapsto a_k$ and $\langle \alpha | E(z) | \alpha \rangle = \frac{u}{1-q}$; $\langle \alpha | F(z) | \alpha \rangle = \frac{u^{-1}}{1-q^{-1}}$.
Then this representation is isomorphic to~$\mathcal{F}_u$.
\end{Proposition}

\begin{proof} Consider the current \[T(z) =\exp \! \left( - \sum_{k>0} \frac{q^{-k/2}-q^{k/2}}{k} a_{-k} z^k \right) E(z) \ \exp \! \left( - \sum_{k<0}\frac{q^{-k/2}-q^{k/2}}{k} a_{-k} z^k \right).\]

It is easy to verify that $[a_k, T(z)] =0$. Since $F_{\alpha}$ is irreducible, $T(z)= f(z)$ for some formal power series $f(z)$ with $\mathbb{C}$-coefficients. On the other hand, $f(z) = \langle \alpha | E(z) | \alpha \rangle = \frac{u}{1-q} $. This implies~\eqref{eq:E}. The proof of~\eqref{eq:F} is analogous.
\end{proof}

\begin{Proposition} \label{E_l}
Denote $E_l(z) = E_{l, k} z^{-k}$. The action of $E_l(z)$ on Fock representation $\mathcal{F}_u$ is given by the following formula
\begin{gather} \label{eq: El in Fock mod}
E_l(z) \rightarrow \frac{u^l}{1-q^l} \exp \! \left( \sum_{k>0} \frac{q^{-kl/2}-q^{kl/2}}{k} a_{-k} z^k \right) \exp \! \left( \sum_{k<0} \frac{q^{-kl/2}-q^{kl/2}}{k} a_{-k} z^k \right).
\end{gather}
\begin{proof}
The commutation relation \eqref{qDc} implies that formula \eqref{eq: El in Fock mod} holds up to a pre-exponential factor. Also, we see from \eqref{qDc} that
\begin{gather} \label{eq: OPE proof for El}
E(z) E_l(w) = \frac{q^{-1}w}{z-q^{-1}w} E_{l+1}\big(q^{-1}w\big) - \frac{q^l w }{z-q^l w} E_{l+1}(w) + \mathrm{reg}.
\end{gather}
The factor can be found inductively from \eqref{eq: OPE proof for El}.
\end{proof}
\end{Proposition}

\subsection{Free fermion realization} \label{fermi}

In this section we give another construction for the Fock representation of $\qD$. To do this, let us consider the Clifford algebra, generated by $\psi_i$ and $\psi^*_j$ for $i,j \in \mathbb{Z}$ subject to the relations
\begin{gather*}
\{ \psi_i, \psi_j \} =0, \qquad \{ \psi^*_i, \psi^*_j \} =0, %\label{fermi1} \\
\qquad \{ \psi_i , \psi^*_j \} = \delta_{i+j, 0}. % \label{fermi2}
\end{gather*}
Consider the currents
\begin{gather*}
\psi(z) = \sum_{i} \psi_i z^{-i-1}, \qquad \psi^*(z) = \sum_i \psi^*_i z^{-i}.
\end{gather*}
Consider a module $F^\psi$ with a cyclic vector $\n$ and relation
\begin{gather*}
\psi_i \n = 0 \quad \text{for} \ i \geqslant l, \qquad \psi^*_j \n = 0 \quad \text{for} \ j > -l.
\end{gather*}
The module $F^\psi$ is independent of $l$. The isomorphism can be seen from the formulas $\psi^*_{-l} \n = | l+1 \rangle$ and $\psi_{l-1} \n = | l-1 \rangle$.
Let us define the $l$-dependent normal ordered product (to be compatible with $\n$) by the following formulas
\begin{gather}
: \! \psi_i \psi^*_j \! :_{(l)} = - \psi^*_j \psi_i \qquad \text{for } i \geqslant l, \label{norm1} \\
: \! \psi_i \psi^*_j \! :_{(l)} = \psi_i \psi^*_j \qquad \text{for } i < l. \label{norm2}
\end{gather}

\begin{Proposition} \label{prop: level 1 Fermion}
The following formulas determine an action of $\qD$ on $F^{\psi}$:
\begin{gather}
c \mapsto 1, \quad c' \mapsto 0, \quad H_k \mapsto \sum_{i+j=k} \psi_i \psi_j^*, \label{eq:centerFermi} \\
E(z) \mapsto \frac{q^{l} u }{1-q} + u q^{ - 1/2 }z : \! \psi\big( q^{-1/2} z\big) \psi^* \big( q^{1/2} z\big) \! :_{(l)} = u q^{-1/2} z \psi\big( q^{-1/2} z\big) \psi^* \big( q^{1/2} z\big), \label{ferm:E} \\
F(z) \mapsto \frac{q^{-l} u^{-1} }{1-q^{-1}} + u^{-1} q^{1/2} z : \! \psi\big( q^{1/2} z\big) \psi^* \big(q^{-1/2} z\big)\!:_{(l)}= u^{-1} q^{1/2} z \psi\big( q^{1/2} z\big) \psi^* \big(q^{-1/2} z\big). \!\!\!\label{ferm:F}
\end{gather}
\end{Proposition}

Let us denote this representation by $\mathcal{M}_u$.

\begin{Remark}The Products $\psi\big( q^{-1/2} z\big) \psi^* \big( q^{1/2} z\big)$ and $\psi\big( q^{1/2} z\big) \psi^* \big(q^{-1/2} z\big)$ from formulas \eqref{ferm:E}--\eqref{ferm:F} are not normally ordered (see Appendix \ref{nonorm} for a formal definition and some other technical details on the \emph{regular product}). In particular, this reformulation implies that $\mathcal{M}_u$ does not depend on~$l$.
\end{Remark}
